\documentclass[11pt,a4paper]{article}

\usepackage[margin=2.2cm]{geometry}
\usepackage[T1]{fontenc}
\usepackage[utf8]{inputenc}
\usepackage{lmodern}
\usepackage{microtype}
\usepackage{amsmath}
\usepackage{amssymb}
\usepackage{booktabs}
\usepackage{enumitem}
\usepackage{hyperref}
\usepackage{xcolor}

\hypersetup{
  colorlinks=true,
  linkcolor=black,
  urlcolor=blue,
  citecolor=black
}

\title{Deterministic RF Field Sensing with Multi-Day Multi-Participant Wi-Fi CSI\\
System Design and Experimental Protocol}
\author{Internal Engineering Paper}
\date{February 2026}

\begin{document}
\maketitle

\begin{abstract}
Radio-frequency (RF) device-free sensing has matured from proof-of-concept localization prototypes into a broad field covering detection, tracking, and human activity recognition using Wi-Fi and other RF signals. Existing surveys on RF-based device-free localization and RF human sensing summarize a rich literature but also reveal a lack of rigorously controlled, multi-day, multi-participant datasets with explicit RF field modeling and environment documentation \cite{denis2019dfl,shahbazian2023human}. At the same time, integrated sensing and communication (ISAC) and emerging WLAN sensing standards such as IEEE 802.11bf indicate that deterministic and reproducible RF sensing will become a first-class system capability \cite{wang2022isac,du2023bf}. This paper specifies a research-grade deterministic RF field sensing system based on Wi-Fi Channel State Information (CSI) and complementary Bluetooth Low Energy (BLE) RSSI, designed for multi-day, multi-participant data collection in a single room. We define the RF and environment model, dual-radio architecture, excitation and control-plane separation, ground-truth protocol, and dataset schema so that the resulting data can support cross-day generalization studies, multi-person inference, and reproducible ML benchmarks.
\end{abstract}

\section{Introduction}

Radio-frequency device-free localization (DFL) and sensing use changes in the wireless channel to detect, track, and identify human presence without requiring tags or devices on the users. A decade of work has produced a wide range of techniques for RF-based detection, tracking, and identification, with Wi-Fi CSI emerging as a primary sensing primitive \cite{denis2019dfl}. In parallel, RF-based human occupancy and activity detection has evolved into a distinct subfield, driven by applications in building energy management, security, and e-health \cite{shahbazian2023human}.

Deep learning has further accelerated this trend, enabling end-to-end learning of features from large-scale indoor IoT data streams and supporting complex applications such as activity recognition and smart home automation \cite{abdelbasset2021indoor}. At the systems level, the emergence of integrated sensing and communication (ISAC) and ongoing standardization efforts in IEEE 802.11bf formalize Wi-Fi-based sensing as a first-class functionality in WLANs \cite{wang2022isac,du2023bf}.

Despite this progress, many experimental Wi-Fi CSI datasets remain small, short-lived, or weakly documented in terms of environment, RF configuration, and temporal variability. Multi-day, multi-participant experiments are often reported, but the environmental and RF-field constraints that make such datasets reproducible and suitable for cross-day generalization studies are rarely specified in engineering detail \cite{shahbazian2023human}. This paper addresses that gap by defining a deterministic RF field sensing system and protocol designed from the ground up for multi-day CSI-based sensing in a single room, with explicit environment modeling and strict RF invariants.

\section{Background and Related Work}

\subsection{RF-Based Device-Free Localization and Human Sensing}

RF-based DFL techniques use variations in channel measurements (RSSI, CSI, or other PHY-layer metrics) to detect and localize human subjects without tags. A comprehensive survey categorizes DFL research into detection, tracking, and identification, highlighting the important role of CSI in improving the spatial resolution and robustness of these systems \cite{denis2019dfl}. Recent surveys on RF-based human sensing further cover occupancy detection, activity recognition, and vital-sign monitoring and emphasize the need for robust, privacy-preserving sensing modalities integrated into the built environment \cite{shahbazian2023human}.

These surveys demonstrate that device-free Wi-Fi CSI sensing can, in principle, support rich human-centric applications, but they also expose limitations in experimental design: many datasets are constrained to single-day captures, limited participant diversity, or loosely controlled environments \cite{denis2019dfl,shahbazian2023human}.

\subsection{Deep Learning for Indoor IoT Sensing}

Deep learning has been widely applied to indoor IoT scenarios, including positioning, tracking, and activity recognition, often operating on large volumes of sensor and RF data generated in smart environments \cite{abdelbasset2021indoor}. CSI-based representations, in particular, are well-suited for convolutional and recurrent architectures due to their structured subcarrier-time-amplitude format. The surveyed work shows that DL methods can effectively learn discriminative features from RF measurements but generally assume that training and test data share similar environmental and temporal conditions \cite{abdelbasset2021indoor}.

Our system is designed to deliberately introduce and record multi-day, multi-participant variability, so that models can be trained and evaluated under realistic covariate shifts, in line with the longer-term goals identified in DL-based indoor IoT sensing \cite{abdelbasset2021indoor}.

\subsection{Indoor Localization, Navigation, and RF Infrastructure}

A parallel line of work studies indoor localization and navigation using RF fingerprints and geometric methods, including RSSI, CSI, and fine time measurement (FTM), in both 2D and 3D settings \cite{sesyuk2022indoor}. Extensive surveys on indoor positioning highlight the increasing demand for GNSS-like functionality indoors and the reliance on smartphone-class hardware as the primary client platform \cite{retscher2022indoor}. Hybrid approaches such as Differential Wi-Fi (DWi-Fi) fuse trilateration and fingerprinting to improve performance and robustness across indoor environments \cite{retscher2017dwifi}.

While these works are not primarily focused on device-free sensing, they provide important context: robust RF sensing must coexist with communication and localization functions in complex indoor deployment scenarios \cite{sesyuk2022indoor,retscher2017dwifi,retscher2022indoor}. Our design aligns with this perspective by isolating sensing and control planes while remaining compatible with commodity WLAN infrastructure.

\subsection{Integrated Sensing and Communication and IEEE 802.11bf}

ISAC aims to unify sensing and communication functions within a single RF system, sharing spectrum, hardware, and signal processing pipelines \cite{wang2022isac}. A recent survey discusses enabling techniques, applications, and datasets for ISAC systems and identifies indoor human sensing as a prominent use case, particularly in smart buildings and IoT deployments \cite{wang2022isac}. In parallel, the IEEE 802.11bf amendment under development standardizes WLAN sensing features, including new use cases, KPIs, and protocol elements tailored to sensing tasks \cite{du2023bf}.

Our system is designed as a pre-802.11bf, research-oriented platform that mimics ISAC-like determinism and repeatability in a controlled lab environment. By explicitly modeling the RF field, environment, and hardware configuration, we aim to generate datasets that could inform future IEEE 802.11bf-compliant designs and ISAC systems \cite{wang2022isac,du2023bf}.

\section{System Goals and Design Principles}

\subsection{Research-Grade Determinism}

The primary goal is to build a deterministic RF field sensing system suitable for multi-day, multi-participant experiments in a single room. Determinism here means:

\begin{itemize}[leftmargin=*]
\item Fixed RF parameters (frequency, bandwidth, MCS, transmit power) within each run.
\item Stable device placement and antenna orientation across days.
\item Explicitly recorded environmental conditions (geometry, layout, doors/windows, HVAC).
\item Controlled excitation traffic with predictable packet timing.
\end{itemize}

This aligns with the call for reproducible and well-characterized experimental setups in RF DFL and human sensing surveys \cite{denis2019dfl,shahbazian2023human}.

\subsection{Multi-Day, Multi-Participant Variability}

Unlike many single-session datasets, our design explicitly targets multi-day, multi-participant variability:

\begin{itemize}[leftmargin=*]
\item Repeated runs over at least $D \geq 3$ days.
\item Multiple sessions per day with varying time-of-day to capture temporal environmental changes.
\item Multiple participants (e.g., 5--20 subjects) with different body types and movement patterns.
\item Configurable number of simultaneous occupants (1..$N$) in the same room.
\end{itemize}

This directly supports the generalization and robustness questions highlighted in RF human sensing and DL-based indoor IoT literature \cite{shahbazian2023human,abdelbasset2021indoor}.

\subsection{Minimal Intrusiveness and Privacy Preservation}

In line with RF-based occupancy and activity detection goals, the system does not require cameras or wearable tags. Ground truth is obtained via scripted protocols, manual annotation, or optional privacy-preserving sensors, consistent with privacy-centric motivations in RF human sensing \cite{shahbazian2023human}.

\section{RF Field and Environment Model}

\subsection{Channel Representation}

We model the Wi-Fi channel for a given link as a complex matrix:
\[
H(f,t) \in \mathbb{C}^{N_{sc} \times N_{rx} \times N_{tx}}
\]
where $N_{sc}$ is the number of subcarriers, $N_{rx}$ and $N_{tx}$ are the number of receive and transmit antennas, $f$ denotes subcarrier frequency, and $t$ denotes time (packet index).

The environment induces a baseline channel $H_0(f)$ in the empty-room state; human presence and movement yield a perturbation $\Delta H(f,t) = H(f,t) - H_0(f)$, which is the primary sensing signal.

\subsection{Room Geometry and Materials}

The room is modeled as a static reflector field characterized by:

\begin{itemize}[leftmargin=*]
\item Dimensions: length $L$, width $W$, height $H$ (in meters).
\item Wall, floor, and ceiling materials (e.g., concrete, drywall, glass).
\item Locations of doors and windows.
\item Furniture layout, including large reflective or absorptive objects.
\end{itemize}

Accurate documentation of these properties is important for both interpretability and cross-room comparison, in line with the emphasis on environment-aware indoor localization systems \cite{sesyuk2022indoor,retscher2022indoor}.

\subsection{Environmental State Variables}

For each run, we log environmental state:

\begin{itemize}[leftmargin=*]
\item Air temperature (°C) and humidity (\%).
\item HVAC state (on/off, mode).
\item Door and window state (open/closed).
\item Time of day.
\item Detected external Wi-Fi networks (basic channel occupancy scan).
\end{itemize}

These variables capture slow drifts and boundary condition changes that can impact channel statistics and are relevant to multi-day generalization studies \cite{shahbazian2023human}.

\section{Hardware and Network Architecture}

\subsection{Receiver Node}

We assume a primary receiver node with:

\begin{itemize}[leftmargin=*]
\item Raspberry Pi 5 as host.
\item Intel AX210 Wi-Fi 6E card for CSI capture (2x2 MIMO).
\item Onboard Wi-Fi/BLE combo (e.g., CYW43455) for BLE scanning and/or control-plane Wi-Fi.
\end{itemize}

CPU and radio temperatures are monitored to detect thermal drift that could influence hardware-level phase stability.

\subsection{Dual-Radio Isolation}

To avoid contamination of the sensing signal with control-plane traffic, we logically separate:

\begin{itemize}[leftmargin=*]
\item \textbf{Sensing radio}: AX210 in monitor or CSI capture mode on a fixed channel and bandwidth.
\item \textbf{Control radio}: Onboard Wi-Fi in a different band (e.g., 2.4\,GHz) or Ethernet for SSH and data transfer.
\end{itemize}

This dual-radio approach aligns with the need to share RF infrastructure between sensing and communication while controlling their mutual interference, as highlighted in ISAC discussions \cite{wang2022isac}.

\subsection{Excitation Sources}

CSI capture requires consistent packet excitation:

\begin{itemize}[leftmargin=*]
\item A fixed access point (AP) or dedicated transmitter on the same channel.
\item Controlled traffic (e.g., UDP stream or CSI-specific injection tool) at a fixed packet rate (e.g., 100--200\,Hz).
\item Optional multiple links (e.g., AP-to-AX210, or multiple APs) for spatial diversity.
\end{itemize}

Excitation should be independent of the control plane to maintain determinism.

\section{Multi-Day Multi-Participant Protocol}

\subsection{Participant Model}

Participants are treated as dynamic scatterers:

\begin{itemize}[leftmargin=*]
\item Multiple individuals with varying height, build, and clothing.
\item Per-run configuration of 0, 1, or $N>1$ occupants.
\item Predefined movement patterns and zones to ensure coverage of the RF field.
\end{itemize}

Multi-person occupancy scenarios are particularly relevant to real-world human sensing applications reviewed in RF occupancy surveys \cite{shahbazian2023human}.

\subsection{Daily Schedule}

A typical multi-day protocol:

\begin{itemize}[leftmargin=*]
\item $D \geq 3$ days.
\item $S \geq 3$ sessions per day (morning, afternoon, evening).
\item Each session includes:
  \begin{itemize}
  \item Empty-room baseline (5--10\,min).
  \item Scripted activities for single and multiple participants.
  \item Free-form segments (natural movement).
  \end{itemize}
\end{itemize}

Session timing is randomized to avoid tight coupling between activity types and time-of-day.

\subsection{Activities and Zones}

We define a set of core activities:

\begin{itemize}[leftmargin=*]
\item Still standing or sitting.
\item Slow and normal walking along predefined paths.
\item Start/stop transitions.
\item Multi-person coordinated and uncoordinated movement.
\end{itemize}

The room is partitioned into zones (e.g., near transmitter, center, near receiver, corners), and each activity is executed in different zones across sessions to sample spatial variability.

\subsection{Ground Truth Annotation}

Ground truth consists of labeled time intervals:

\begin{itemize}[leftmargin=*]
\item Occupancy count (0,1,2,\dots).
\item Activity label (from a controlled vocabulary).
\item Optional zone label and confidence score.
\end{itemize}

Annotation sources may include manual logs, synchronized button presses, or privacy-preserving sensors. This reflects common practices in RF human sensing experiments while respecting privacy considerations \cite{shahbazian2023human}.

\section{Signal Processing and Feature Extraction}

\subsection{CSI Preprocessing}

For each packet, we obtain $H(f,t)$ and compute:

\begin{itemize}[leftmargin=*]
\item Amplitude $A(f,t) = |H(f,t)|$.
\item Sanitized phase after removing linear and common phase components, where feasible.
\item Per-chain RSSI if available.
\end{itemize}

Subcarriers corresponding to guard bands or unstable pilots are excluded.

\subsection{Temporal Windowing}

Two window scales serve different tasks:

\begin{itemize}[leftmargin=*]
\item Short windows (10--50\,ms) for motion and micro-Doppler features.
\item Long windows (0.5--2\,s) for occupancy-level estimation.
\end{itemize}

This multi-scale design is consistent with reported use of windowed features for RF-based activity recognition and human sensing \cite{shahbazian2023human,abdelbasset2021indoor}.

\subsection{Feature Sets}

For each window we compute:

\begin{itemize}[leftmargin=*]
\item Subcarrier amplitude statistics (mean, variance, median absolute deviation).
\item Temporal derivatives and their energy.
\item Short-time Fourier transform features (Doppler-like signatures).
\item Packet rate and loss indicators.
\item BLE RSSI summary statistics (mean, variance, packet rate) per advertising device.
\end{itemize}

These features can be consumed directly by ML models or used as input to deeper architectures as in DL-based indoor sensing systems \cite{abdelbasset2021indoor}.

\section{Dataset Structure and ML Interface}

\subsection{Run-Level Organization}

Each run is stored as:

\begin{itemize}[leftmargin=*]
\item \texttt{meta.json}: environment, hardware, RF configuration, timing.
\item \texttt{wifi/}: raw CSI, packet metadata, windowed features.
\item \texttt{ble/}: raw advertisements and aggregated features.
\item \texttt{labels/}: ground truth intervals and derived per-window labels.
\item \texttt{qa/}: quality metrics, packet rate plots, drift reports.
\end{itemize}

This structure is designed for direct ingestion into ML pipelines and aligns with the need for high-quality datasets in ISAC and RF sensing research \cite{wang2022isac}.

\subsection{Train/Validation/Test Protocols}

We define several evaluation regimes:

\begin{itemize}[leftmargin=*]
\item \textbf{Within-day}: train and test on the same day (session split).
\item \textbf{Cross-day}: train on day $D_1$, test on $D_2,\dots$.
\item \textbf{Cross-participant}: leave-one-participant-out evaluation.
\item \textbf{Cross-occupancy}: training on single-person, testing on multi-person scenarios.
\end{itemize}

These regimes reflect generalization challenges identified in RF human sensing and indoor IoT DL work \cite{shahbazian2023human,abdelbasset2021indoor}.

\section{Discussion and Future Directions}

Our specification aligns with long-term trends identified in RF DFL, RF-based human sensing, and ISAC: sensing will be a built-in capability of indoor RF infrastructure, and datasets must support cross-day, cross-participant, and cross-environment generalization \cite{denis2019dfl,shahbazian2023human,wang2022isac,du2023bf}. By explicitly modeling the RF field, environment, and temporal variability, this system aims to provide a foundation for such research.

Future work includes:

\begin{itemize}[leftmargin=*]
\item Extending to multi-room and multi-floor deployments, informed by 3D indoor localization insights \cite{sesyuk2022indoor}.
\item Integrating smartphone-collected data channels to bridge with user-centric navigation systems \cite{retscher2022indoor}.
\item Exploring few-shot calibration and transfer learning across rooms and buildings.
\end{itemize}

\section{Conclusion}

We have presented a research-grade deterministic RF field sensing system and experimental protocol for multi-day, multi-participant Wi-Fi CSI and BLE-based human sensing in a single room. The design is grounded in the state of the art in RF-based device-free localization, RF human sensing, deep learning for indoor IoT, and emerging ISAC and WLAN sensing standards \cite{denis2019dfl,shahbazian2023human,abdelbasset2021indoor,wang2022isac,du2023bf}. By enforcing strict RF invariants, detailed environmental logging, dual-radio isolation, and carefully structured multi-day protocols, the resulting datasets are intended to support robust evaluation and development of models that can operate reliably in real indoor environments.

\begin{thebibliography}{9}

\bibitem{denis2019dfl}
S.~Denis, R.~Berkvens, and M.~Weyn,
``A Survey on Detection, Tracking and Identification in Radio Frequency-Based Device-Free Localization,''
\textit{Sensors}, 2019.
Available via Dimensions: \url{https://app.dimensions.ai/details/publication/pub.1123145036}.

\bibitem{abdelbasset2021indoor}
M.~Abdel-Basset, V.~Chang, H.~Hawash, R.~K. Chakrabortty, and M.~Ryan,
``Deep learning approaches for human-centered IoT applications in smart indoor environments: a contemporary survey,''
\textit{Annals of Operations Research}, 2021.
Available via Dimensions: \url{https://app.dimensions.ai/details/publication/pub.1139558226}.

\bibitem{du2023bf}
R.~Du, H.~Hua, H.~Xie, X.~Song, Z.~Lyu, M.~Hu, Narengerile, Y.~Xin, S.~McCann, M.~Montemurro, T.~X.~Han, and J.~Xu,
``An Overview on IEEE 802.11bf: WLAN Sensing,''
\textit{arXiv preprint}, 2023.
Available via Dimensions: \url{https://app.dimensions.ai/details/publication/pub.1165325343}.

\bibitem{sesyuk2022indoor}
A.~Sesyuk, S.~Ioannou, and M.~Raspopoulos,
``A Survey of 3D Indoor Localization Systems and Technologies,''
\textit{Sensors}, 2022.
Available via Dimensions: \url{https://app.dimensions.ai/details/publication/pub.1153341341}.

\bibitem{retscher2017dwifi}
G.~Retscher,
``FUSION OF LOCATION FINGERPRINTING AND TRILATERATION BASED ON THE EXAMPLE OF DIFFERENTIAL WI-FI POSITIONING,''
\textit{ISPRS Annals of the Photogrammetry Remote Sensing and Spatial Information Sciences}, 2017.
Available via Dimensions: \url{https://app.dimensions.ai/details/publication/pub.1091774366}.

\bibitem{retscher2022indoor}
G.~Retscher,
``Indoor Navigation—User Requirements, State-of-the-Art and Developments for Smartphone Localization,''
\textit{Geomatics}, 2022.
Available via Dimensions: \url{https://app.dimensions.ai/details/publication/pub.1154046398}.

\bibitem{shahbazian2023human}
R.~Shahbazian and I.~Trubitsyna,
``Human Sensing by Using Radio Frequency Signals: A Survey on Occupancy and Activity Detection,''
\textit{IEEE Access}, 2023.
Available via Dimensions: \url{https://app.dimensions.ai/details/publication/pub.1157479341}.

\bibitem{wang2022isac}
J.~Wang, N.~Varshney, C.~Gentile, S.~Blandino, J.~Chuang, and N.~Golmie,
``Integrated Sensing and Communication: Enabling Techniques, Applications, Tools and Data Sets, Standardization, and Future Directions,''
\textit{IEEE Internet of Things Journal}, 2022.
Available via Dimensions: \url{https://app.dimensions.ai/details/publication/pub.1149475587}.

\end{thebibliography}

\end{document}
